% ============================================================================
% Fichier  : partie5/P05_C18_agilite_devops_cloud.tex
% Titre    : Agilité, DevOps et Cloud : Gouverner à Grande Vitesse
% Auteur   : MINKA MI NGUIDJOÏ Thierry Emmanuel
% Version  : 1.0 -- octobre 2026
% ============================================================================
% Sources : Manifeste agile (2001) ; Forsgren, Humble et Kim (2018) ; NIST
% SP 800-145 (2011) ; ISO/IEC 27017 et 27018 ; loi n° 2010/012 (texte
% anglais officiel) ; loi n° 2024/017 (présentations publiées). Le tableau
% comparatif africain ne cite que l'intitulé et l'année des lois ; aucun
% article n'est cité.
% ============================================================================

\chapter{Agilité, DevOps et Cloud : Gouverner à Grande Vitesse}
\label{chap:agile-cloud}
\chaptermark{Agilité, DevOps et cloud}

\epigraph{%
    « La vitesse n'est pas l'ennemie du contrôle.\\
    C'est le contrôle manuel qui est l'ennemi de la vitesse. »%
}{Principe des équipes DevOps}

\niveauChapitre{Expert}{%
    Maîtrise des Chapitres~\ref{chap:itil-changements} (gestion des
    changements) et~\ref{chap:itil-niveaux} (niveaux de service et
    fournisseurs).%
}

\begin{boxobjectifs}
\begin{itemize}
    \item Concilier contrôle et agilité (DevOps, gestion automatisée des
          changements).
    \item Gouverner l'informatique en nuage (responsabilité partagée,
          réversibilité, localisation des données).
    \item Évaluer un contrat d'informatique en nuage au regard du droit
          camerounais.
    \item Comparer les approches de gouvernance numérique de plusieurs pays
          africains.
\end{itemize}
\end{boxobjectifs}

\bigskip

% ============================================================================
\section*{Question directrice : peut-on gouverner ce qui change chaque jour~?}
% ============================================================================

La nouvelle équipe mobile de la BANQUE SANAGA déploie une version de
l'application chaque semaine, et voudrait passer à un déploiement par
jour. Le comité des changements, qui se réunit le mardi, ne peut pas
suivre. À Bafoussam, NOUN PAY n'a pas un seul serveur à elle~: tout son
système d'information est loué, à la minute, à un fournisseur de nuage
étranger.

Les Parties~\ref{part:cobit} à~\ref{part:pratique} ont construit une
gouvernance faite de comités, de procédures et d'autorisations. La
Partie~\ref{part:avance}, qui s'ouvre ici, la confronte aux pratiques
qui en bousculent les rythmes et les frontières. Ce chapitre traite de
l'agilité et de DevOps (sections~\ref{sec:c18-agile}
et~\ref{sec:c18-devops}), de l'informatique en nuage
(section~\ref{sec:c18-cloud}), du contrat de nuage au regard du droit
camerounais (section~\ref{sec:c18-contrat}), puis des approches
africaines de la gouvernance numérique (section~\ref{sec:c18-afrique}).

\separateur

% ============================================================================
\section{L'Agilité}
\label{sec:c18-agile}
% ============================================================================

\subsection{Le manifeste et ses conséquences}

Le \emph{Manifeste pour le développement agile de logiciels}, publié en
2001 par dix-sept praticiens, pose quatre préférences~: les individus et
leurs interactions plutôt que les processus et les outils, des logiciels
opérationnels plutôt qu'une documentation exhaustive, la collaboration
avec les clients plutôt que la négociation contractuelle, l'adaptation au
changement plutôt que le suivi d'un plan~\cite{agile2001manifeste}. Il
précise que les éléments de droite gardent leur valeur, mais que ceux de
gauche en ont davantage.

\begin{boxalert}
\textbf{Violation typique.} Une équipe invoque l'agilité pour ne rien
documenter, ne rien faire autoriser et ne rendre compte à personne. Lors
d'un incident, personne ne sait quelle version est en production ni qui
l'a validée.

\smallskip
\textbf{Règle~:} l'agilité change la \emph{forme} du contrôle, pas son
existence. Le manifeste ne supprime ni la documentation, ni les contrats,
ni les plans~: il en change le poids relatif.
\end{boxalert}

\subsection{Ce qui change pour la gouvernance}

\begin{tableaucours}{L{3.2cm}L{4.4cm}L{4.6cm}}{%
    \tableauHeader{Sujet} &
    \tableauHeader{Approche traditionnelle} &
    \tableauHeader{Approche agile}
}
Décision d'investissement & Un projet, un budget, une date de fin & Un
produit, une équipe stable, un financement par période \\
Pilotage & Respect du plan, du budget, des délais & Valeur livrée à chaque
itération \\
Exigences & Définies au départ & Priorisées en continu par un responsable
de produit \\
Contrôle & Jalons et validations formelles & Démonstrations régulières,
critères de « fini » \\
\end{tableaucours}
\vspace{-0.8em}
\begin{center}{\small\itshape Tableau~18.1. Gouvernance traditionnelle et
gouvernance agile}\end{center}

La gouvernance ne disparaît pas~: elle se déplace. Le comité stratégique
n'approuve plus un cahier des charges, il fixe les objectifs d'un produit,
l'enveloppe d'une équipe et les indicateurs de valeur
(Chapitre~\ref{chap:alignement}), puis revoit ces décisions à intervalles
réguliers.

% ============================================================================
\section{DevOps}
\label{sec:c18-devops}
% ============================================================================

\subsection{Rapprocher le développement et l'exploitation}

\begin{boxdef}
\termecle{DevOps}\index{DevOps}~: ensemble de pratiques culturelles et
techniques qui rapproche les équipes de développement et d'exploitation,
pour livrer des changements fréquents, fiables et réversibles, grâce à
l'automatisation de la construction, des tests et du déploiement.

\medskip
\termecle{Chaîne de déploiement}\index{chaîne de déploiement}
(\textit{pipeline})~: suite automatisée des étapes par lesquelles un
changement de code passe, de sa validation par un pair à sa mise en
production.
\end{boxdef}

\subsection{Mesurer la performance de la livraison}

Les travaux de Forsgren, Humble et Kim~\cite{forsgren2018accelerate} ont
popularisé quatre indicateurs de la performance de la livraison
logicielle. Ils ont une propriété précieuse pour la gouvernance~: ils
mesurent à la fois la vitesse \emph{et} la stabilité.

\begin{tableaucours}{L{3.8cm}L{4.4cm}L{4.0cm}}{%
    \tableauHeader{Indicateur} &
    \tableauHeader{Mesure} &
    \tableauHeader{Dimension}
}
Fréquence de déploiement & Nombre de mises en production par période &
Vitesse \\
Délai de livraison des changements & Temps entre la validation du code et
sa mise en production & Vitesse \\
Taux d'échec des changements & Part des déploiements qui provoquent un
incident ou un retour arrière & Stabilité \\
Délai de rétablissement & Temps pour rétablir le service après un échec &
Stabilité \\
\end{tableaucours}
\vspace{-0.8em}
\begin{center}{\small\itshape Tableau~18.2. Les quatre indicateurs de la
performance de livraison}\end{center}

Le taux d'échec des changements et le délai de rétablissement sont les
indicateurs de BAI06 et de DSS02 (Chapitres~\ref{chap:itil-changements}
et~\ref{chap:itil-incidents})~: DevOps ne remplace pas la gestion des
changements, il en mesure le résultat.

\subsection{Déplacer le contrôle dans la chaîne}

\begin{tableaucours}{L{4.0cm}L{8.2cm}}{%
    \tableauHeader{Contrôle traditionnel} &
    \tableauHeader{Équivalent dans une chaîne DevOps}
}
Autorisation par le comité des changements & Changement standard
pré-autorisé, si la chaîne valide toutes les étapes~; comité réservé aux
changements hors norme \\
Séparation des tâches entre développeur et exploitant & Revue obligatoire
du code par un pair~; le développeur ne peut pas fusionner seul son code \\
Tests de recette manuels & Tests automatisés obligatoires, avec seuil de
couverture bloquant \\
Plan de retour arrière & Déploiement progressif et retour arrière
automatique en cas d'erreur \\
Procès-verbal d'autorisation & Journal inaltérable de la chaîne (qui, quoi,
quand, avec quels résultats de tests) \\
\end{tableaucours}
\vspace{-0.8em}
\begin{center}{\small\itshape Tableau~18.3. Les contrôles de la gestion des
changements dans une chaîne DevOps}\end{center}

\begin{boxscenario}{Opération SANAGA. L'équipe mobile}
\textbf{Demande.} L'équipe mobile veut déployer chaque jour, sans passer
par le comité des changements.

\medskip
\textbf{Décision du comité stratégique.} Les déploiements de l'application
mobile sont reconnus comme \emph{changements standard} à quatre conditions,
vérifiées par la chaîne~: revue du code par un pair, tests automatisés
couvrant au moins 90~\% des fonctions de paiement, déploiement progressif
(5~\% des clients pendant une heure), retour arrière automatique. Tout
changement touchant l'interface avec le système central reste un
changement normal, soumis au comité.

\medskip
\textbf{Contrôle.} L'auditrice interne examine chaque trimestre un
échantillon de déploiements dans le journal de la chaîne~: elle teste
désormais un contrôle \emph{automatisé}, et non plus une signature.
\end{boxscenario}

\begin{boxmetier}{Auditeur SI}
\textbf{Auditer une chaîne DevOps}
\begin{enumerate}
    \item Obtenir la description de la chaîne et des contrôles qu'elle
          impose, puis vérifier qu'ils sont bloquants (un test échoué
          arrête le déploiement).
    \item Vérifier qui peut modifier la chaîne elle-même~: c'est le
          nouveau point de contrôle critique.
    \item Rapprocher les déploiements en production du journal de la
          chaîne~: un déploiement hors chaîne est un changement non
          autorisé.
    \item Suivre les quatre indicateurs de performance de livraison, et
          notamment le taux d'échec.
\end{enumerate}
\end{boxmetier}

% ============================================================================
\section{L'Informatique en Nuage}
\label{sec:c18-cloud}
% ============================================================================

\subsection{Définitions}

La définition de référence est celle du NIST~\cite{nist2011cloud}~:
l'informatique en nuage permet un accès à la demande, par le réseau, à un
ensemble partagé de ressources informatiques configurables, rapidement
mises à disposition et libérées avec un effort de gestion minimal. Elle
présente cinq caractéristiques~: libre-service à la demande, large accès
par le réseau, mutualisation des ressources, élasticité rapide, service
mesuré.

\begin{tableaucours}{L{2.8cm}L{4.8cm}L{4.6cm}}{%
    \tableauHeader{Modèle} &
    \tableauHeader{Ce que fournit le prestataire} &
    \tableauHeader{Exemple}
}
IaaS (infrastructure) & Serveurs, stockage, réseau virtuels & L'hébergement
des serveurs de NOUN PAY \\
PaaS (plateforme) & Environnement d'exécution, bases de données gérées &
La base de données gérée du modèle de notation \\
SaaS (logiciel) & Application complète & La messagerie de la BANQUE SANAGA \\
\end{tableaucours}
\vspace{-0.8em}
\begin{center}{\small\itshape Tableau~18.4. Les trois modèles de service du
NIST}\end{center}

\subsection{La responsabilité partagée}

Dans le nuage, la responsabilité de la sécurité et de l'exploitation est
\emph{partagée} entre le client et le fournisseur, selon le modèle de
service. Quel que soit le modèle, \textbf{le client reste responsable de
ses données, de ses utilisateurs et de ses droits d'accès}.

\begin{tableaucours}{L{4.0cm}ccc}{%
    \tableauHeader{Domaine} &
    \tableauHeader{IaaS} &
    \tableauHeader{PaaS} &
    \tableauHeader{SaaS}
}
Données et leur classification & Client & Client & Client \\
Identités et droits d'accès & Client & Client & Client \\
Application & Client & Client & Fournisseur \\
Système d'exploitation, intergiciels & Client & Fournisseur & Fournisseur \\
Infrastructure physique, réseau & Fournisseur & Fournisseur & Fournisseur \\
\end{tableaucours}
\vspace{-0.8em}
\begin{center}{\small\itshape Tableau~18.5. Le modèle de responsabilité
partagée (simplifié)}\end{center}

\begin{boxalert}
\textbf{Violation typique.} « Nos données sont chez un grand fournisseur
mondial, elles sont donc en sécurité. » Le stockage de NOUN PAY était
configuré en accès public par erreur~; le fournisseur, dont la plateforme
fonctionnait parfaitement, n'y était pour rien.

\smallskip
\textbf{Règle~:} la sécurité du nuage est celle du fournisseur~; la
sécurité \emph{dans} le nuage est celle du client.
\end{boxalert}

\subsection{Gouverner le nuage}

Le nuage déplace quatre risques que la gouvernance doit arbitrer
explicitement~:
\begin{itemize}
    \item \textbf{la dépendance}~: changer de fournisseur peut être long et
          coûteux, surtout en PaaS~;
    \item \textbf{la localisation}~: les données peuvent être stockées
          hors du territoire, ce que le droit encadre
          (section~\ref{sec:c18-contrat})~;
    \item \textbf{l'informatique fantôme}~: n'importe quel service peut
          souscrire un abonnement avec une carte bancaire, hors de tout
          contrôle~;
    \item \textbf{le coût}~: payé à l'usage, il dérive vite sans pilotage.
\end{itemize}
Deux normes aident à traiter la sécurité~: l'ISO/IEC~27017, qui donne des
mesures de sécurité propres aux services en nuage, et l'ISO/IEC~27018, qui
traite de la protection des données personnelles dans le nuage public.

% ============================================================================
\section{Le Contrat de Nuage et le Droit Camerounais}
\label{sec:c18-contrat}
% ============================================================================

\begin{boxjuris}
\textbf{\loicam{n°~2010/012}, articles~26 et~29} (version anglaise
officielle, extraits)~:

\medskip
\textit{« [Operators of information systems] shall have standardized
systems enabling them to at all times identify, assess, process or manage
any risk relating to the security of the information systems of the
services provided directly or indirectly. »} (art.~26, al.~1)

\medskip
\textit{« Operators of information systems shall be bound to conserve the
connection and traffic data of their information systems for a period of
10 (ten) years. »} (art.~29, al.~1)

\medskip
\textbf{\loicam{n°~2024/017}} (d'après les présentations publiées). La loi
encadre strictement le transfert des données personnelles hors du
territoire national (Chapitre~\ref{chap:droit}).

\medskip
\textbf{Portée pour un contrat de nuage.} Une organisation qui confie son
système d'information à un fournisseur de nuage reste tenue de maîtriser
les risques de sécurité des services fournis « indirectement », de
conserver ses données de connexion pendant dix ans, et de faire examiner
tout transfert de données personnelles hors du Cameroun au regard de la
loi n°~2024/017. Le contrat type d'un fournisseur mondial n'a pas été
écrit pour ces obligations~: il faut en vérifier chaque point.
\end{boxjuris}
% A VERIFIER : conditions exactes des transferts hors du territoire dans le
% texte publié de la loi n° 2024/017 et ses décrets d'application.

\begin{boxprocedure}
\textbf{La grille d'évaluation d'un contrat de nuage.} Avant de signer, ou
lors de l'audit, chaque point est noté conforme, partiellement conforme
ou non conforme, avec la clause correspondante~:
\begin{itemize}
    \item localisation des données et des sauvegardes, et liste des pays~;
    \item conditions de transfert hors du Cameroun~;
    \item conservation et accès aux journaux pendant la durée légale~;
    \item droit d'audit, ou rapports d'audit indépendants et leur
          périmètre~;
    \item notification des incidents de sécurité, dans un délai fixé~;
    \item engagements de service et pénalités~;
    \item sous-traitants du fournisseur et information du client~;
    \item réversibilité~: formats, délais, assistance, coûts~;
    \item droit applicable et juridiction compétente.
\end{itemize}
\end{boxprocedure}

\begin{boxscenario}{Opération NOUN. Deux contrats comparés}
\begin{tableaucours}{L{4.2cm}L{3.8cm}L{3.8cm}}{%
    \tableauHeader{Point de la grille} &
    \tableauHeader{Contrat type} &
    \tableauHeader{Contrat d'entreprise}
}
Localisation des données & Non conforme (non choisie) & Conforme (région
choisie, liste des pays) \\
Droit d'audit & Non conforme (renvoi aux certifications) & Partiel
(rapports d'audit indépendants) \\
Notification des incidents & Partiel (« dans les meilleurs délais ») &
Conforme (72~heures) \\
Réversibilité & Non conforme & Partiel (exports, sans assistance) \\
Journaux & Partiel (90~jours) & Partiel (configurable, coût de stockage) \\
\end{tableaucours}

\medskip
\textbf{Analyse} (confiance moyenne). Le contrat d'entreprise, pour un
surcoût de 18~millions FCFA par an (Chapitre~\ref{chap:cas-integrateurs}),
corrige les points les plus graves. Il ne règle pas tout~: la conservation
des journaux pendant dix ans suppose un archivage complémentaire, et la
réversibilité doit être préparée par NOUN PAY elle-même (formats ouverts,
exercice de restitution annuel).
\end{boxscenario}

\begin{boxmetier}{RSSI}
\textbf{Gouverner les abonnements en nuage}
\begin{enumerate}
    \item Tenir l'inventaire des services en nuage utilisés, y compris ceux
          payés par carte bancaire par les métiers.
    \item Classer les données selon leur sensibilité, et interdire les
          plus sensibles sur les services non évalués.
    \item Vérifier la configuration des services (accès publics, chiffrement,
          comptes à privilèges) au moins chaque mois.
    \item Présenter chaque année au comité des risques la carte des
          dépendances aux fournisseurs de nuage.
\end{enumerate}
\end{boxmetier}

% ============================================================================
\section{La Gouvernance Numérique en Afrique}
\label{sec:c18-afrique}
% ============================================================================

\subsection{Un cadre continental}

L'Union africaine a adopté en 2020 une \emph{Stratégie de transformation
numérique pour l'Afrique 2020--2030}, puis en 2022 un \emph{Cadre
stratégique sur les données}, qui encourage l'harmonisation des règles de
circulation des données sur le continent. La Convention de Malabo, entrée
en vigueur en 2023, en constitue le volet juridique
(Chapitre~\ref{chap:droit}). L'initiative Smart Africa, lancée en 2013,
réunit des États autour de projets numériques communs.

\subsection{Des lois nationales de protection des données}

\begin{tableaucours}{L{3.0cm}L{6.2cm}L{3.0cm}}{%
    \tableauHeader{Pays} &
    \tableauHeader{Texte} &
    \tableauHeader{Année}
}
Sénégal & Loi sur la protection des données à caractère personnel & 2008 \\
Maroc & Loi n°~09-08 relative à la protection des personnes physiques à
l'égard du traitement des données personnelles & 2009 \\
Ghana & Data Protection Act & 2012 \\
Côte d'Ivoire & Loi relative à la protection des données à caractère
personnel & 2013 \\
Kenya & Data Protection Act & 2019 \\
Rwanda & Loi sur la protection des données personnelles et de la vie
privée & 2021 \\
Nigeria & Nigeria Data Protection Act & 2023 \\
Cameroun & Loi n°~2024/017 & 2024 \\
\end{tableaucours}
\vspace{-0.8em}
\begin{center}{\small\itshape Tableau~18.6. Quelques lois africaines de
protection des données (intitulés abrégés)}\end{center}

\subsection{Trois enseignements}

\begin{enumerate}
    \item \textbf{La localisation des données est un choix politique.}
          Certains États privilégient la libre circulation pour attirer les
          services, d'autres la conservation sur le territoire pour des
          raisons de souveraineté. Une organisation présente dans plusieurs
          pays doit en tenir un registre (Chapitre~\ref{chap:droit}).
    \item \textbf{Les infrastructures nationales se développent.} Plusieurs
          États ont investi dans des centres de données nationaux, qui
          offrent une alternative au nuage étranger~; leur niveau de
          service et de sécurité doit être évalué comme celui de tout
          fournisseur.
    \item \textbf{Le droit précède souvent les autorités.} Une loi n'a
          d'effet que si l'autorité qu'elle crée dispose de moyens.
          L'efficacité réelle d'un régime se mesure à ses décisions
          publiées, encore rares (Chapitre~\ref{chap:defaillances}).
\end{enumerate}

\begin{boiteRemarque}{Souveraineté numérique ou isolement~?}
Imposer que toutes les données restent au Cameroun protège contre
certains risques (accès par des autorités étrangères, coupure de
connectivité internationale), mais en crée d'autres~: dépendance à un
nombre réduit de fournisseurs locaux, coûts plus élevés, moindre
résilience si un seul centre de données héberge tout. La bonne réponse
de gouvernance n'est ni l'interdiction ni l'indifférence, mais la
\emph{classification}~: localiser ce qui doit l'être, et choisir
librement pour le reste, contrat par contrat.
\end{boiteRemarque}

% ============================================================================
\separateur
% ============================================================================

\begin{boxresume}
\textbf{Ce qu'il faut retenir de ce chapitre~:}
\begin{itemize}
    \item \textbf{L'agilité change la forme du contrôle, pas son
          existence}~: la gouvernance finance des produits et des équipes,
          et revoit ses décisions à intervalles réguliers.
    \item \textbf{DevOps déplace le contrôle dans la chaîne}~: revue par un
          pair, tests bloquants, déploiement progressif, retour arrière
          automatique, journal inaltérable. Le point critique devient la
          chaîne elle-même.
    \item \textbf{Quatre indicateurs} mesurent ensemble la vitesse et la
          stabilité de la livraison.
    \item \textbf{Dans le nuage, la responsabilité est partagée}~: le client
          reste responsable de ses données, de ses utilisateurs et de ses
          accès.
    \item \textbf{Le contrat de nuage s'évalue au regard du droit
          camerounais}~: risques des services fournis indirectement,
          conservation des journaux pendant dix ans, transferts hors du
          territoire.
    \item \textbf{L'Afrique se dote de règles}~: stratégie continentale,
          lois nationales de protection des données, centres de données
          nationaux. La réponse de gouvernance est la classification des
          données.
\end{itemize}
\end{boxresume}

\bigskip

\begin{boxexo}[title={\ding{45}\quad Exercice~18.1. Qui est responsable~? \niveaubase}]
Pour chacune des situations suivantes, indiquez si la responsabilité
incombe au client ou au fournisseur de nuage, selon le Tableau~18.5.
\begin{enumerate}[(a)]
    \item Un ancien salarié de NOUN PAY se connecte encore à la console
          d'administration.
    \item Une panne électrique arrête le centre de données du fournisseur.
    \item La messagerie en SaaS de la BANQUE SANAGA présente une faille
          logicielle.
    \item Une base de données gérée (PaaS) est accessible depuis Internet
          sans mot de passe.
\end{enumerate}
\end{boxexo}

\begin{boxexo}[title={\ding{45}\quad Exercice~18.2. Calculer les indicateurs de livraison \niveaumoyen}]
En un mois, l'équipe mobile de la BANQUE SANAGA a réalisé 22~déploiements.
Trois ont provoqué un incident~; les délais de rétablissement ont été de
20~minutes, 45~minutes et 3~heures. Le délai moyen entre la validation du
code et la mise en production est de 2~jours.
\begin{enumerate}[(a)]
    \item Calculez le taux d'échec des changements et le délai moyen de
          rétablissement.
    \item L'incident de 3~heures a touché le paiement. Qu'en déduisez-vous
          sur le périmètre des changements standard~?
    \item Proposez une cible pour chacun des quatre indicateurs.
\end{enumerate}
\end{boxexo}

\begin{boxexo}[title={\ding{45}\quad Exercice~18.3. Migrer l'ANPP vers le nuage~? \niveauexpert}]
Un fournisseur propose à l'ANPP d'héberger son système de paiement dans
son nuage, hors du Cameroun, pour un coût annuel inférieur de 30~\% à celui
de son centre de données actuel.
\begin{enumerate}[(a)]
    \item Classez les données de l'agence selon leur sensibilité.
    \item Appliquez la grille d'évaluation du contrat de nuage à l'offre,
          en indiquant les informations qui vous manquent.
    \item Identifiez les obligations légales en jeu, en citant uniquement
          les articles vérifiés, et signalez ce qui reste à vérifier.
    \item Rédigez la recommandation au conseil d'administration, en
          présentant au moins deux options.
\end{enumerate}
\end{boxexo}

\bigskip

\begin{boxinfo}
\textbf{Transition.} Le nuage concentre les données~; l'intelligence
artificielle les exploite pour décider. Le Chapitre~\ref{chap:donnees-ia}
traite de la gouvernance des données et de l'intelligence artificielle, à
partir du cas de la notation de crédit de NOUN PAY.
\end{boxinfo}
